\section{来源审计与课程主线：为什么 Audio 是 Transformer 的压力测试}

本讲不是把 NLP Transformer 简单搬到声音上，而是追问更基础的设计问题：连续、高采样率、跨多个时间尺度的声音，应该先被表示成什么，再交给什么预测目标？Prateek Verma 用三篇工作串起 raw waveform synthesis、self-supervised representation learning 与 large-scale audio understanding。三条线看似分别对应生成、表征和分类，实际共享同一条工程链：\term{representation} 决定序列长度与信息损失，\term{objective} 决定模型保留什么，\term{metric} 决定我们究竟证明了什么。

官方课程页将这场课列在 2021-11-29，Stanford Online 到 2022-07-18 才上传录像。视频对应的 Audio Transformers 论文在课堂时是 2021 年 v1；arXiv 后来于 2025 年更新了 v2。本讲义只用 v1 解释课堂实验，不能把四年后的修订倒灌成讲者当时已有的证据。视频未公开独立 slide PDF，因此从 1080p 官方录像的 1,449 个采样点中人工筛选 47 个 teaching states，并用 981 条官方人工字幕恢复老师的口头解释。

\lecturefigure{slide-01-title-understanding-to-synthesis.jpg}{课程标题把主线概括为：从 language modeling 到 understanding，再到 synthesis。}{Stanford Online 官方视频 00:00:24--00:00:47。}

标题页中的三个动词不是营销分类。\emph{Synthesis} 要生成高频波形；\emph{understanding} 要把长信号压成对标签有用的表示；\emph{language modeling} 则要求先把连续声音变成可预测的离散符号。读后文时应不断问：当前系统处理的是 sample、spectrum、embedding 还是 code index？如果这一层没说清楚，所谓“Audio Transformer”就只是一个模糊名称。

\teachervoice{Verma 一开场把这场课称为 ``buffet style''：他不会证明一种统一架构必胜，而是让三篇论文展示多种可组合思路。讲义因此按共同问题组织，而不按论文摘要机械分章。}

\subsection{三篇论文怎样组成一条设计链}

第一篇工作在 8-bit raw samples 上做 next-sample prediction，比较 WaveNet 与 causal Transformer；第二篇先把连续音频变成 codebook tokens，再比较 contrastive learning 与 next-code prediction；第三篇直接让 Transformer 从 raw waveform 学分类表示，并在中间 embedding 上加入 pooling 或 Haar wavelet。三篇工作的连接点如下：

\lecturefigure{slide-02-talk-roadmap.jpg}{课堂路线图：表示基础、raw synthesis、离散 token 表征学习、raw-waveform understanding。}{Stanford Online 官方视频 00:01:08--00:01:45。}

\begin{knowledgebox}{读图：不要把三条路线混成“同一个 Transformer”}
Raw synthesis 的输出是下一个量化 sample；representation learning 的输出可以是 code prediction 或下游 linear probe；audio understanding 的输出是多标签 sound-event scores。它们共享 attention，但输入粒度、监督信号、序列长度和评价指标都不同。正确比较单位不是模型名字，而是“表示--目标--指标”三元组。
\end{knowledgebox}

\begin{warningbox}{课堂时间边界}
讲者在 2021 年把 Transformer 称为跨领域的“当前风潮”，并用 “Adieu WaveNet?”、 “Adieu Convolutions” 作为挑衅性标题。这些标题只能解释当时三组受控实验，不能推出 2026 年所有语音、音乐、codec 或生成系统都应放弃卷积、状态空间模型或扩散模型。
\end{warningbox}

\subsection{证据应该按什么强度阅读}

后文既有代数事实，也有 benchmark table、滤波器可视化和讲者的未来判断。为了不把它们写成同一强度的结论，本讲采用四级证据：结构上能否实现、受控实验是否改善、内部表示是否可解释、结论能否外推到感知质量或其他任务。任何一级都不能自动替代下一级。

\begin{table}[H]
\centering
\small
\caption{本讲常见证据与能够支持的结论。}
\begin{tabularx}{\textwidth}{L{0.21\textwidth}L{0.29\textwidth}X}
\toprule
证据 & 例子 & 合理结论 \\
\midrule
结构与复杂度 & $n=f_sT$、attention 为 $O(n^2)$ & 说明为什么 raw audio 需要局部窗口、压缩或稀疏结构。 \\
受控指标 & 同数据下 top-5 accuracy 或 mAP & 说明该设置中的预测/分类性能，不等于主观音质或普适架构优势。 \\
表示探针 & Linear probe、学习滤波器的频响分析 & 说明表示携带可用信息或出现类似经典信号处理的结构，不证明单个 neuron 有唯一语义。 \\
外推与判断 & “更多数据会缩小 gap”“Transformer 不是终点” & 是研究假设或方向，应与实测结果分开。 \\
\bottomrule
\end{tabularx}
\end{table}

\subsection{本章小结}

本讲的对象不是单一模型，而是音频建模的设计空间。课堂日期、上传日期和论文修订日期必须分开；三篇工作的结果也必须在各自输入表示、目标函数与指标范围内解释。

